\section{Related Work}
\label{sec:related_work}

% ── Instructions ─────────────────────────────────────────────────────────
% Organise into subsections by theme. Each subsection: survey existing work,
% identify limitation, link to your approach. Aim for 40–60 references total.
% ─────────────────────────────────────────────────────────────────────────

\subsection{Graph-Based SLAM}

Graph-based SLAM~\cite{grisetti2010tutorial} represents the robot trajectory
as a factor graph whose nodes are poses and whose edges encode odometry or
loop-closure constraints.
Optimisation with Gauss–Newton or Levenberg–Marquardt~\cite{ceres-solver}
minimises the total squared residual.
A representative implementation~\cite{macenski2021slam} extends this
with lifelong and multi-session mapping, operating in synchronous or
asynchronous scan-matching modes.
A key limitation of purely LiDAR-based graphs is that drift accumulates in
featureless corridors where loop-closure opportunities are rare.

\subsection{Infrastructure Camera Localisation}

Infrastructure-mounted cameras offer a complementary sensing modality for
indoor localization because they observe the robot from viewpoints that are
independent of the robot's onboard motion.
However, as emphasized in YOWO~\cite{yang2025yowo}, reliable registration of
ceiling-mounted cameras to the scene layout remains the key bottleneck.
The related literature highlighted by YOWO falls into three main groups.

The first group uses SLAM or visual localization to reconstruct an indoor scene
and then localize a query camera against that scene model.
Representative methods such as InLoc~\cite{taira2018inloc} and hierarchical
visual localization~\cite{sarlin2019coarse} achieve strong performance in
indoor visual localization, but they rely heavily on static scene features and
can degrade under viewpoint disparity, repeated structures, and ambiguous
textures.
These challenges become more severe when matching ceiling-mounted cameras
against egocentric observations collected from very different perspectives.

The second group estimates relative poses between infrastructure cameras by
tracking mobile keypoints, often from a human moving through the scene.
Examples include extrinsic camera calibration from a moving person
~\cite{lee2022extrinsic} and online marker-free extrinsic calibration using
person keypoints~\cite{patzold2022online}.
These methods are robust to cross-view ambiguity because they rely on dynamic
agents rather than static scene appearance, but they typically recover only
relative camera geometry and do not automatically align the resulting camera
poses to a globally consistent scene map.

The third group studies collaboration between egocentric and external cameras.
Works such as EgoHumans~\cite{khirodkar2023egohumans} show that combining
egocentric and externally mounted views improves downstream perception tasks,
but camera registration is generally handled beforehand as a separate step.
YOWO is important because it explicitly connects these three lines of research
and argues for a unified formulation that jointly maps the scene and registers
ceiling-mounted cameras.

Our work shares YOWO's motivation that post-hoc registration is fragile, but it
targets a different estimation setting: instead of RGB-D scene reconstruction
with collaborative ego and overhead video, we focus on 2D LiDAR pose-graph SLAM
and use fixed cameras as static anchor sources for reducing drift in the robot
trajectory.

\subsection{Constraint Injection in SLAM}

Suppressing cumulative drift by injecting external constraints into the SLAM
optimization graph is a well-established idea.
Two common sources of such constraints are semantic landmarks and
infrastructure-based ranging anchors.
For example, fusing UWB range measurements with short-range 2D LiDAR can add
global distance constraints that reduce the position drift and map divergence of
standalone 2D SLAM~\cite{liu2022uwb2dlidar}.
However, UWB-based approaches remain sensitive to indoor non-line-of-sight and
multipath effects, and their performance often depends on a sufficiently dense,
pre-installed anchor infrastructure~\cite{liu2022uwb_nlos,jiang2025gpsdenied}.

Semantic or vision-aided SLAM methods inject higher-level environmental cues to
improve loop closure and map consistency.
For instance, low-cost LiDAR-vision fusion has been used to build 2.5D maps and
stabilize localization by combining complementary geometric and visual
information~\cite{jiang2019lidarvision}.
Yet these methods can become computationally demanding and may degrade in
environments with repetitive appearance, dynamic objects, or weakly distinctive
textures.

Structurally, our method is closest to SLAM systems that bound drift by adding
absolute reference measurements to a pose graph.
In outdoor mapping, Global Positioning System (GPS) observations play this role
by constraining the long-term drift of LiDAR odometry.
Surveys of GPS-denied LiDAR SLAM make clear, however, that such global fixes are
unavailable in indoor facilities and deep urban canyons, where signal
attenuation or occlusion renders satellite positioning unreliable or unusable
~\cite{jiang2025gpsdenied}.
Our method follows the same design principle of augmenting the graph with
external absolute references, but replaces GPS or UWB range anchors with
camera-derived robot observations from fixed infrastructure cameras.
